\chapter{Discussion and Conclusion}
\label{ch:conclusion}

%-------------------------------------------------------

%%%%%%%%%%%%%%%%%%%%%%%%%%%%%%%%%%%%%%%%%%%%%%%%%%%%%
\section{Discussion}
\label{sec:discussion}
%%%%%%%%%%%%%%%%%%%%%%%%%%%%%%%%%%%%%%%%%%%%%%%%%%%%%

% TODO: interpret the results. Suggested points:
%   - what the histogram comparison reveals about the point operation
%   - where information is lost (saturation / clipping) and why
%   - performance: does the display keep up with the source frame rate

As a software engineer, I like to split every project into a backend part and a frontend part; separating them from the start lets me reason about the two sets of technologies independently. By the backend I mean the calculations that change the image's settings --- vector multiplication and so on. To actually see the result, rather than just debugging from the console, a frontend is needed. After some research, Tkinter was chosen because it is very simple to use: for now the interface only needs to be functional, and Tkinter gets the job done while keeping the code simple.


%%%%%%%%%%%%%%%%%%%%%%%%%%%%%%%%%%%%%%%%%%%%%%%%%%%%%
\section{Challenges and Limitations}
\label{sec:limitations}
%%%%%%%%%%%%%%%%%%%%%%%%%%%%%%%%%%%%%%%%%%%%%%%%%%%%%

% TODO: be honest and specific. Problems encountered, how each was
% resolved, and what the program still does not handle.
The main challenge was finding and organizing the functions needed to perform the operations. It ended up being a matter of identifying the right OpenCV (\texttt{cv2}) function for each task: \texttt{cv2.VideoCapture} and \texttt{cv2.VideoWriter} to read and write the video files, \texttt{cv2.calcHist} to build the histograms, \texttt{cv2.convertScaleAbs} to apply the brightness/contrast point operation with built-in clipping to $[0,255]$, and \texttt{cv2.resize}/\texttt{cv2.cvtColor} to prepare each frame for display. After that, the next challenge was connecting these results to the UI. OpenCV frames are NumPy arrays in BGR channel order, while Tkinter widgets can only display \texttt{PhotoImage} objects, so each frame had to be converted to RGB and wrapped with \texttt{PIL.ImageTk.PhotoImage} before it could be shown on a \texttt{Label} --- and a reference to that image had to be kept alive manually (Section \ref{sec:task1}), since Tkinter does not hold its own reference and would otherwise garbage-collect the image and leave the panel blank.


%%%%%%%%%%%%%%%%%%%%%%%%%%%%%%%%%%%%%%%%%%%%%%%%%%%%%
\section{Conclusion}
\label{sec:conclusion}
%%%%%%%%%%%%%%%%%%%%%%%%%%%%%%%%%%%%%%%%%%%%%%%%%%%%%

Brightness and contrast are point operations: they depend only on the input pixel to calculate the output. In this assignment, we showed how this works by adding sliders that change the video's brightness and contrast in real time, and observed how the histogram changes according to the values assigned to gain and bias.
